\begin{document}

\section{Abstract}

We model an agentic runtime as $\mathcal{A}=\langle\mathcal{S},\mathcal{H},\mathcal{X}\rangle$: task Skills specify what may be done, a Harness compiles probabilistic intent into governed execution, and Scaffold instances provide isolated physical capacity. The central contribution is a bounded, falsifiable separability conjecture. Within a declared operating region, it predicts that logical capability and physical capacity can vary independently when six causal conditions hold: typed closure, complete mediation, effect non-interference, shared-state isolation, resource invariance, and scheduler independence. The Harness exposes a contract boundary whose interface semantics comprise the admission decision, activated path, authorized effects, binding constraints, and postconditions. Deterministic gates and trace identity make violations observable, auditable, and attributable; they are not asserted to cause independence. We provide protocols that estimate semantic invariance, capacity response, and recoupling interactions rather than synthetic results. We also describe architecture-compatible mechanisms for context partitioning, external-data access, an Intermediate Relation between data and output registries, locality-aware execution, planning-time dry-runs, and a train-then-freeze Skill lifecycle with monitoring. These mechanisms are design patterns, secondary hypotheses, or open questions, not logical consequences of the conjecture. The intended result is a research agenda and an implementable vocabulary for agentic runtimes whose growth remains observable, governable, and reversible.

\section{Introduction}
\label{sec:introduction}

The practical problem in an agentic runtime is not simply how to make a large language model (LLM) call more tools. It is how to let a system gain capabilities and capacity without making either dimension opaque. A runtime may add a database Skill, a code-analysis Skill, and a procurement Skill while serving the same number of requests. It may also add workers, sandboxes, accelerators, or regional pools while exposing exactly the same capabilities. The first operation expands a logical action space; the second expands a physical execution space. Combining them in one undifferentiated agent loop hides the distinction that operations teams most need to see.

Consider a common scaling maneuver: every tool description, schema, policy, and example is loaded into every request so that the model can decide what to use. This appears flexible, but it turns capability registration into prompt growth. Unrelated Skills compete for attention, control metadata crosses into request context, and the effect of adding one capability is difficult to localize. A second maneuver adds more workers behind the same loop. Throughput may rise, yet safety and replay do not improve because the runtime still lacks a contract that explains what was planned, which checks ran, and why a particular worker was authorized. More capacity amplifies an execution model; it does not repair it.

This paper asks a narrower question than general agent intelligence: \textit{what runtime boundary permits logical capability scaling and physical throughput scaling to evolve independently while remaining manageable?} Our answer is a testable conjecture about the Harness contract. Skill does not directly depend on Scaffold; Scaffold does not understand task-specific Skill semantics. Their dependencies meet only through typed Harness contracts. The conjecture is neither a theorem nor a universal independence result. It predicts separability inside a declared operating region only under typed closure, complete mediation, effect non-interference, shared-state isolation, resource invariance, and scheduler independence. Departures from these conditions must be detectable, auditable, and attributable. Deterministic gates and trace identity supply that observability; they do not generate separability by themselves.

The architecture is organized around a compact model:

\begin{equation}
\mathcal{A}=\langle\mathcal{S},\mathcal{H},\mathcal{X}\rangle,
\end{equation}

where $\mathcal{S}$ is a versioned set of Skills, $\mathcal{H}$ is the Harness that compiles and governs runs, and $\mathcal{X}$ is a pool of Scaffold instances. A Skill declares a capability contract. The Harness selects Skills, constructs an execution graph, applies policy and budget checks, binds graph nodes to Scaffold instances, and records evidence. A Scaffold supplies resources and isolation without interpreting the business meaning of the Skill.

\begin{figure}[t]
\centering
\caption{Dual scaling model. Logical extension adds or revises Skills through the Harness contract; physical extension adds Scaffold capacity through the same binding boundary. The architecture makes capability growth and capacity growth distinct interventions whose separability can be tested inside a declared operating region.}
\label{fig:dual-scaling}
\end{figure}

Figure 1 establishes the paper's organizing distinction. Horizontal growth changes what the runtime can do; vertical growth changes how much governed work it can execute. The Harness is not merely glue between them. It is the control surface that makes both forms of growth inspectable.

\subsection{Contributions}

This paper makes three primary contributions.

\begin{enumerate}
\item We state the Testable Separability Conjecture: within a declared operating region, typed closure, complete mediation, effect non-interference, shared-state isolation, resource invariance, and scheduler independence are hypothesized to permit Skill capability and Scaffold capacity to vary independently.
\item We define the Harness as a typed capability-to-capacity binding boundary. Its interface semantics are the admission decision, activated path, authorized effects, binding constraints, and postconditions. Deterministic gates and trace identities make those semantics inspectable and violations attributable.
\item We provide falsification methodology rather than synthetic results. A randomized capability-count $\times$ Scaffold-count experiment estimates semantic invariance, capacity response, and recoupling interaction; additional protocols test control-plane leakage, composed-path safety, locality and dry-run economics, Skill lifecycle stability, vector-valued training admission, reconstructive reflection, and registry independence.
\end{enumerate}

The paper also develops architecture-compatible mechanisms: high-cohesion context partitioning, derivation-closure orchestration, external-data contracts, parallelism and locality policies, planning-time dry-runs, an Intermediate Relation, and a train-freeze-monitor Skill lifecycle. The Harness boundary makes these patterns implementable and testable, but does not make them necessary or guarantee their performance.

\section{Origins of the Extensibility Framework}
\label{sec:origins}

\subsection{The Cloud Application Analogy}

The three-axis decoupling proposed here is not invented from scratch. Modern cloud applications have long operated with an implicit three-layer separation that maps directly onto our runtime objects. Virtual machines provide isolated execution capacity; software frameworks (application servers, message queues, API gateways) provide the runtime contract layer between business logic and infrastructure; and front-end applications compose user-facing capabilities on top. In this analogy, adding a new web page is a logical extension that need not require provisioning more virtual machines; adding more virtual machines need not change which pages exist. These examples motivate a stable contract boundary, but do not establish separability for a probabilistic agentic runtime.

The Skill-Harness-Scaffold model adapts this pattern to an agentic runtime. Skills correspond to front-end application code: they declare what the user-facing capability is, not how the machine runs. Scaffolds correspond to virtual machines: they provide isolated execution envelopes whose count scales throughput. The Harness corresponds to the framework runtime: it compiles declarative intent into executable units, enforces cross-cutting policy, and binds logical requests to physical resources. The mapping motivates the testable hypothesis that independent scaling, fault isolation, and contract-based composition can hold inside a declared operating region when typed closure, complete mediation, effect non-interference, shared-state isolation, resource invariance, and scheduler independence hold. The shared capability-count $\times$ Scaffold-count factorial and its diagnostic arms must determine whether that prediction survives contact with an implemented system.

\subsection{What Is Genuinely New: Probabilistic Calls in Every Layer}

The analogy also reveals the precise boundary of novelty. In a traditional cloud application, every layer executes deterministic code. In an agentic runtime, every layer may invoke an LLM: the Skills layer generates natural-language plans, the Harness performs contract translation and tool synthesis using model calls, and even the Scaffold layer's serving engine processes probabilistic inference. This creates a fundamentally new problem: the contract boundary between layers must mediate not just deterministic interfaces but probabilistic ones. A Skill declaration may be ambiguous; a Harness translation may vary across runs; a Scaffold's inference output is inherently stochastic.

This is the gap that existing work does not close. The policy-driven runtime layer of Zhang et al. \cite{zhang2026policy} identifies the ``seam'' between agent frameworks (which understand semantics) and serving engines (which process events), and inserts an intermediate layer for cross-cutting policies. But its four primitives, observe, score, predict, and act, assume deterministic coordination of serving-level events. Our Harness contract goes further: it must mediate probabilistic intent (from Skills), probabilistic translation (within the Harness itself), and probabilistic execution (on Scaffolds), while producing deterministic gates and replayable traces.
The core technical question is: under what causal conditions can semantic behavior remain invariant as physical capacity changes, and capacity response remain stable as logical capability changes? Section~5.3 states six candidate conditions and an experiment that can count against them.

\subsection{Control-Plane and Data-Plane Orthogonality}

A second source of the extensibility thesis comes from distributed systems practice. The separation of control plane (low-frequency decisions, configuration, routing) from data plane (high-frequency payload processing) is a widely used design pattern in SDN, Kubernetes, database query execution, and service meshes. The pattern motivates independent optimization of control complexity, per-request processing cost, and data-plane throughput, but its presence alone does not establish independence.

In an agentic runtime, this separation provides a second motivation for the separability conjecture. Logical scaling (adding Skills) primarily enlarges the control-plane decision space by introducing new capability contracts, new selection rules, and new composition invariants, without directly loading more tokens into the data plane. Physical scaling (adding Scaffolds) primarily enlarges the data-plane throughput capacity. Acting on different planes makes the two forms of growth distinct interventions, not structurally orthogonal ones: shared state, resources, effects, and scheduler feedback can still couple them. Moreover, the agentic control plane is not the deterministic configuration distribution of SDN; it is a \textit{probabilistic control plane} whose decisions may be LLM-generated. Deterministic gates and auditable traces make those decisions inspectable, while the six causal conditions and the factorial protocol determine whether semantic invariance and capacity response remain separable.

\subsection{The Data Subsystem as a Third Axis}

A third source emerges from the observation that agents do not only scale in capabilities and capacity; they also scale in the breadth of external data they must understand. An agent that starts with one data source may grow to access dozens: enterprise databases, SaaS platforms, personal files, and public internet sources. If each new source injects its schema, content, and metadata into the request context, data growth becomes a third form of prompt pollution. The data subsystem resolves this by pushing semantic summarization into an independent off-policy loop, an asynchronous background process with its own agents and its own model, that pre-summarizes sources into a queryable index. The main request path only consumes the query interface, not the full source content. This makes data scaling analogous to physical scaling: adding sources expands what is reachable without proportionally increasing per-request cost. Chen et al. \cite{chen2026metadata} empirically demonstrate the tension between pre-built metadata (high precision, low coverage) and on-line discovery (high coverage, low precision); the off-policy loop is designed to occupy the gap they leave open. The temporal knowledge graph architecture of Zep \cite{zep2025} and the 5W1H-guided knowledge graph construction of SoKG \cite{sokg2026}, together with the event-ontology framework of XPEventCore \cite{xpeventcore2024}, inform the structure of the Data Wiki summaries introduced in Section~8.2. The data subsystem in this paper is intentionally scoped to its contract-relevant part: the Intermediate Relation. Its full four-component formalization (on-policy APIs, data-usage skills, lifecycle, off-policy loop) is developed in companion notes and placed outside our scope in Section~11.

\section{Related Work and Novelty Boundary}
\label{sec:related}

\paragraph{Runtime and serving layers.}
The closest work is a policy-driven runtime layer for coordinating agent frameworks and serving engines \cite{zhang2026policy}. That line of work identifies an important systems boundary: policies chosen above the serving engine need runtime support below it. Our focus is complementary. We ask how task capabilities themselves are represented, activated, composed, and bound to isolated execution capacity. The distinction matters because scheduling policy can be correct while the capability path remains untyped or unaudited.

Jagarin \cite{kadaboina2026jagarin} and Helium \cite{wadlom2026helium} explore runtime support for efficient agent execution. Dynamic runtime-graph work \cite{xie2026runtimegraphs} similarly moves beyond static workflow templates by constructing graphs at run time. We adopt dynamic graphs inside the Harness, but use \textit{Scaffold} in a narrower sense: execution and isolation infrastructure such as a sandbox, worker, container, accelerator allocation, or regional pool. It should not be confused with a workflow scaffold or prompt scaffold.

Tool Forge \cite{rao2026toolforge} treats tool production and adaptation as a systems problem. Skill-as-Code in this paper shares the emphasis on lifecycle, but its unit is a validated capability contract rather than only an executable tool. A Skill may include a tool adapter, a planner fragment, schemas, policy tags, tests, and postconditions. Five-plane runtime governance \cite{tallam2026fiveplane} provides a broader governance decomposition; our Harness can be understood as the point at which such governance planes become enforceable on a particular activated path. ClawVM \cite{clawvm2026} places compaction and writeback in a harness layer with deterministic and auditable guarantees, but its contract scope is narrower: it governs memory operations rather than the full capability-to-capacity binding.

\paragraph{Concurrent work (July 2026).}
The work closest to our manageability claim appeared while this paper was in preparation. Soni \cite{soni2026gates} instantiates a discipline we treat as a design requirement: every capability must pass a pre-declared machine-checkable suite, and a fixed invariant set, including the property that no action reaches an effector without a capability token issued by a control ring, is model-checked over the reachable states of a bounded model across successive releases. That work reports governance overhead of roughly 0.021~ms per request. We regard it as the strongest available evidence that deterministic gating need not be expensive, and we adopt its measurement discipline for our own protocols. Its scope, however, is a single runtime's safety core verified over a bounded model. It does not treat logical capability growth and physical throughput growth as separately scalable axes, which is the specific separation this paper studies, and its invariants govern the action-to-effector path rather than the capability-to-capacity binding.

Two further July 2026 results bear directly on our Skill lifecycle. Skillware \cite{fan2026skillware} supplies the software ontology we had left implicit, separating a Skill Artifact from the Skillware Unit that carries its identity and from the Agent Host that activates it, and defining lifecycle continuity, whether identity survives update, rollback, and removal, as a measurable property over a corpus of more than 138{,}000 Skill records. We adopt that vocabulary rather than restate it; Section~8.4 accordingly narrows its claim to the freeze-to-code stage and its role as a deterministic anchor. Scoped verification of evolving agentic context \cite{hsu2026grace} stores a persistent system instruction as a typed semantic graph and validates proposed edits only within the local typed neighborhood of modified nodes, reporting a large reliability gain over a flat-text baseline. This is independent evidence for the principle behind our Intermediate Relation, namely that structure is what makes verification local; the difference is the object being structured, which for that work is the instruction and for us is the coupling point between a data registry and an output registry. Related harness-evolution work makes behavior localization itself the bottleneck \cite{wang2026handbook} or adapts the harness from an experience bank \cite{huang2026memoharness}; both operate inside the maintenance subsystem that Section~11 places outside our scope.

Empirical studies of the Skill ecosystem also inform our assumptions. Large-scale corpora show that Skill descriptions and actual behavior can diverge \cite{zhang2026misalignment,wang2026skillcorpus}, which is the practical obstacle to the operation closure that Section~5.1 requires, and lifecycle-aware threat modeling over real Skills finds risks well beyond execution time \cite{badhe2026skillsecurity}, supporting the path-aware and release-time checks we specify.

Composition safety is especially relevant. Xie et al. show that components benign in isolation may become harmful in composition \cite{xie2026composition}. This observation directly challenges per-tool allowlists as a sufficient runtime defense. Our response is path-aware checking: the object of authorization is the activated Skill path plus its data and resource bindings, not a bag of individually approved Skills.

\paragraph{Production frameworks.}
Production agent frameworks and orchestration engines occupy adjacent regions of the design space but do not expose the typed contract boundary studied here. AutoGen \cite{wu2023autogen} coordinates multi-agent conversations but treats tool registration as prompt-level metadata rather than a versioned contract with deterministic gates. LangGraph \cite{langgraph2024} constructs stateful graph workflows but conflates planning and execution in a single runtime without a separate binding boundary to isolated capacity. CrewAI \cite{crewai2024} assigns roles to agents but provides no path-level invariant checking or replay identity. The OpenAI Assistants API \cite{openai2024assistants} bundles tools, context, and execution into one managed service where the contract between capability and capacity is implicit. Temporal \cite{temporal2024} provides durable execution and replay for workflows but operates at the orchestration layer without agent-specific capability contracts or composition safety. AgentScope \cite{gao2024agentscope} and Agent-as-a-Service \cite{agentasaas2025} offer multi-agent platforms and scalable infrastructure, respectively, but neither separates logical capability scaling from physical throughput scaling as distinct axes governed by a typed contract. Autellix \cite{liu2025autellix} is the closest on the physical axis, modeling agents as general programs for serving throughput, but does not address the logical scaling axis.

Table~1 summarizes where these systems fall relative to five Harness responsibilities. It covers production frameworks and workflow engines only; the concurrent research discussed above is placed by argument rather than by rating, since its objects of governance are not the same.

\begin{table}[t]
\centering
\caption{Comparison of Harness responsibilities across production agent and workflow systems. Ratings record only coverage verifiable in public documentation and cited sources; ``not documented'' does not establish that a capability is absent. The Proposed design row reports intended properties, not empirical results. The table is restricted to systems that register task capabilities and execute them. Concurrent research with a different governance object is positioned in Section~4 rather than scored here.}
\label{tab:system-comparison}
\begin{tabular}{lccccc}
\hline
System & Skill & Contract & Policy & Binding & Replay \\
 & registration & compilation & gates & to capacity & / audit \\
\hline
AutoGen & partial & not documented & not documented & not documented & not documented \\
LangGraph & partial & not documented & not documented & implicit & partial \\
CrewAI & partial & not documented & not documented & not documented & not documented \\
Assistants API & documented & not documented & implicit & implicit & not documented \\
Temporal & documented & partial & documented & documented & documented \\
Helium & not documented & not documented & not documented & partial & not documented \\
Tool Forge & documented & partial & partial & not documented & partial \\
Five-Plane & documented & not documented & documented & not documented & documented \\
\textbf{Proposed design}\textsuperscript{†} & \textbf{proposed} & \textbf{proposed} & \textbf{proposed} & \textbf{proposed} & \textbf{proposed} \\
\hline
\end{tabular}
\end{table}

Prior agent-memory and retrieval-augmented generation (RAG) systems address context persistence and evidence retrieval \cite{packer2023memgpt,lewis2020rag}. Those mechanisms can be hosted behind a Skill contract, but they do not by themselves establish the capability-to-capacity boundary studied here. SkillOpt \cite{wang2026skillopt} demonstrates that skills are the trainable external state of a frozen agent, with optimizer-level discipline applied in text space; our Skill-as-Code lifecycle extends this by adding a freeze-to-code stage that anchors determinism. CodeAct \cite{wang2024codeact} shows that executable code actions improve agent composability; our work positions code not merely as an efficiency gain but as a deterministic anchor in a probabilistic control plane. The novelty claim is therefore deliberately limited: this paper does not propose a new model architecture, scheduler algorithm, memory representation, or universal governance taxonomy. It proposes a contract boundary and a set of conditions and tests for scalable, manageable agentic runtime behavior.

\section{Model and Assumptions}
\label{sec:model}

\subsection{The Three Runtime Objects}

\paragraph{Skill.}
A Skill $s\in\mathcal{S}$ is a versioned declaration of task capability. Its minimum interface is

\begin{equation}
s=\langle g,i,o,p,e,v\rangle,
\end{equation}

where $g$ describes the goal and applicability conditions, $i$ and $o$ are typed input and output schemas, $p$ contains preconditions and policy tags, $e$ declares externally visible effects, and $v$ identifies the validated version. A Skill can include model instructions, tool adapters, graph fragments, examples, and tests, but these assets are implementation details behind the contract. A Skill is \textit{operation-closed} when every external effect it can request is declared and testable through its interface.

Collection-scale evidence suggests that operation closure is not what current practice provides. The SkillMD-138K snapshot released with the Skillware ontology \cite{fan2026skillware} contains 138,133 content-deduplicated \texttt{SKILL.md} records drawn from 20,556 repositories, and two of its reported measurements bear on closure. First, frontmatter is nearly universal (136,380 records, 98.73\%), which the authors read narrowly as evidence that a structured metadata envelope is a dominant visible convention, explicitly not that its fields are validated or that activation semantics agree across agent systems. Second, a lexical detector finds an explicit path token into \texttt{scripts/}, \texttt{references/}, or \texttt{assets/} in 32,069 records (approximately 23.2\%); the authors note that this detector misses implicit dependencies, so the figure is a lower bound on packaged artifacts rather than a census. Even read as a lower bound, it indicates that a large share of the corpus is natural-language text with no declared deterministic execution component. Neither figure speaks directly to the fields operation closure requires. A frontmatter envelope establishes that a name and a description are conventional, which corresponds to the goal field $g$; it does not establish a typed output schema $o$, a declared effect set $e$, or a version identity $v$ that a gate could reject independently. The corpus statistics therefore bound what can be claimed: a metadata envelope is the prevailing convention, whereas the typed and independently rejectable interface this architecture depends on is not visible at collection scale. That gap is the reason the Harness must enforce closure as an admission condition rather than assume it of registered Skills. Whether closure can be retrofitted onto existing Skills at acceptable cost is an empirical question we do not settle here.

\paragraph{Harness.}
The Harness $\mathcal{H}$ is the runtime compiler and governor. Given a request $q$, identity and policy context $u$, and a Skill registry $\mathcal{S}$, it produces a candidate execution graph $D$, validates the graph, binds accepted nodes to Scaffold instances, and emits a trace:

\begin{equation}
(q,u,\mathcal{S}) \xrightarrow{\mathcal{H}} (D,C,b,\tau).
\end{equation}

Here $C$ is the resolved Harness contract, $b$ is a binding from graph nodes to Scaffold instances, and $\tau$ is the audit and replay trace. Selection and graph construction may use an LLM and are therefore probabilistic. Admission, schema validation, policy evaluation, budget enforcement, effect authorization, and trace creation must be deterministic for a fixed contract and policy snapshot.

\paragraph{Scaffold.}
A Scaffold instance $x\in\mathcal{X}$ is a physical execution envelope:

\begin{equation}
x=\langle r,z,l,a\rangle,
\end{equation}

where $r$ describes available resources, $z$ describes isolation and trust zone, $l$ describes data and network locality, and $a$ describes attestation and runtime identity. A Scaffold may be a process sandbox, container, virtual machine, browser session, graphics processing unit (GPU) allocation, or remote execution pool. It exposes resource facts and execution primitives, not task semantics.

\subsection{The Harness Contract}

The resolved contract for one run is

\begin{equation}
C=\langle I,O,G,B,E,T\rangle.
\end{equation}

$I$ and $O$ are the run-level typed inputs and outputs. $G$ is the activated Skill graph. $B$ contains resource, time, token, cost, concurrency, placement, and isolation constraints. $E$ contains authorized effects and evidence obligations. $T$ contains policy, model, registry, data-snapshot, and trace identities needed for replay.
The boundary's interface semantics are defined by five observable components: the admission decision over $C$, the activated path in $G$, the authorized effects in $E$, the binding constraints in $B$, and the postconditions encoded by $O$ and its evidence obligations. This tuple is intentionally operational: each field must be serializable, inspectable, and independently rejectable.

\begin{figure}[t]
\centering
\caption{Anatomy of a resolved Harness contract. A probabilistic plan becomes executable only after schemas, graph structure, budgets, effects, evidence obligations, and trace identities have passed deterministic gates. Manageability resides in the transition from generated intent to a typed, auditable execution unit.}
\label{fig:harness-contract}
\end{figure}

Figure 2 makes the Harness boundary concrete. A useful shorthand is: \textit{probabilistically written, deterministically and auditably executed}. The phrase does not imply that execution outcomes are deterministic; external services and models remain variable. It means that the authorization decision, declared effects, selected versions, and evidence requirements can be reconstructed from the recorded contract.

\subsection{Conditional Decoupling}

The separability conjecture is scoped to a declared operating region $\Omega$. A study of $\Omega$ fixes the workload class; model, tokenizer, and policy versions; external-service contracts; admissible resource classes; and scheduler regime. A change outside those bounds is a change to the operating region, not an isolated $\Delta\mathcal{S}$ or $\Delta\mathcal{X}$ intervention. Within $\Omega$, the conjecture names six causal conditions. They are hypotheses about when semantic behavior and capacity response can vary independently, not a proof that the architecture guarantees independence.

\begin{enumerate}
\item \textbf{Typed closure.} The activated operation and effect surface resolves to versioned typed input, output, and effect declarations; operations or effects outside that surface are rejected before execution.
\item \textbf{Complete mediation.} Every Scaffold, tool, and external-system invocation traverses an instrumented Harness binding point; there is no alternate execution or call channel outside that boundary.
\item \textbf{Effect non-interference.} Effect-labelled data or control flow from one activated path can reach another only through a dependency declared in their contracts; undeclared cross-path read, write, and control channels are absent.
\item \textbf{Shared-state isolation.} Every mutable store is assigned to a declared namespace, snapshot, or synchronization protocol, and cross-partition access is possible only through a declared shared-state dependency.
\item \textbf{Resource invariance.} Scaffold count, topology, and identity have no channel into semantic inputs: across compatible Scaffold arms, model, tokenizer, policy, numerical-mode, prompt and tool payload, timeout and retry, and sampling configurations are unchanged except for declared binding fields and resource-only telemetry.
\item \textbf{Scheduler independence.} Scheduler-labelled signals are confined to declared ordering and resource-dependency fields and have no undeclared flow into admission, path-selection, effect-authorization, binding-constraint, or postcondition logic.
\end{enumerate}

Before outcome collection, each study must preregister a separate mechanism-level acceptance criterion and its instrumentation for every condition: (i) typed closure passes only if every activated operation resolves to the declared versioned types and all preregistered undeclared-operation and undeclared-effect probes fail closed; (ii) complete mediation passes only if every observed invocation carries a valid Harness binding and negative probes expose zero unbound call paths; (iii) effect non-interference passes only if information-flow traces and cross-path probes expose zero undeclared read, write, or control edges; (iv) shared-state isolation passes only if every mutable access maps to its declared namespace, snapshot, or synchronization protocol and probes expose zero undeclared cross-partition accesses; (v) resource invariance passes only if configuration and payload hashes for the semantic-input fields listed above are equal across compatible Scaffold arms and resource-labelled taint reaches none of those fields; and (vi) scheduler independence passes only if scheduler-perturbation traces expose zero scheduler-labelled flows outside the declared ordering and resource-dependency fields. These condition criteria are evaluated without semantic task scores, interface-semantic outcomes, throughput, latency, or factorial interactions. An outcome leaving its margin does not itself fail a condition, and passing all six condition checks does not establish separability.

Under these candidate conditions, adding $s_{new}$ to $\mathcal{S}$ is predicted to leave capacity response unchanged for existing workloads, while adding compatible $x_{new}$ to $\mathcal{X}$ is predicted to preserve semantic behavior and the five interface semantics. The Harness contract is the candidate stable interface. Shared mutable state, dynamic policy, nondeterministic external effects, resource-sensitive model behavior, cross-Skill interference, and scheduler feedback can re-couple the axes even when declared interfaces remain unchanged; the operating region and evaluation protocol must therefore expose these factors rather than assume them away.

\begin{proposition}[Testable Separability Conjecture]
\label{prop:p1}
Within a declared operating region $\Omega$, if typed closure, complete mediation, effect non-interference, shared-state isolation, resource invariance, and scheduler independence meet their independently preregistered mechanism-level acceptance criteria, then interventions on Skill capability count and compatible Scaffold count are predicted to be separable: Scaffold changes preserve preregistered semantic behavior and the five interface semantics, while Skill changes preserve the capacity-response function for existing workloads.
\end{proposition}

The six conditions identify concrete recoupling channels without asserting necessity. Lack of typed closure permits undeclared behavior; incomplete mediation permits direct binding; interfering effects and shared state transmit changes between paths; resource observations can enter semantic inputs; and scheduler signals can enter task-control logic. A system may compensate for one weakness through a mechanism outside this model, so observing separability after an apparent condition violation would refine the conjecture rather than establish a converse.

Deterministic gates form part of a separate observability requirement. Gates make the admission decision, activated path, authorized effects, binding constraints, and postconditions reproducible for fixed recorded inputs.
Trace identity completes that requirement by linking those decisions to the Skill, model, policy, data snapshot, binding, and effects involved.
Violations become detectable, auditable, and attributable through this requirement, which is not a causal generator of independence. P1 is rejected in the tested operating region if all six independently evaluated condition criteria pass but preregistered semantic equivalence or non-inferiority fails, capacity response departs beyond its margin, or a capability-count $\times$ Scaffold-count interaction exceeds its margin. A non-significant interaction alone does not establish independence.

P1 is the core structural conjecture. P15 (dual-subgoal convergence in Section~8.5), P16 (reconstructive reflection sufficiency), and P17 (Intermediate-Relation decoupling in Section~8.2) address training and data subsystems. Table~\ref{tab:propositions} separates claim type from current evidence status.

\begin{table}[t]
\centering
\caption{Claim types and evidence status. ``Proposed experiment'' denotes a protocol, not a completed study in this paper.}
\label{tab:propositions}
\begin{tabular}{lp{0.14\linewidth}p{0.30\linewidth}p{0.34\linewidth}}
\hline
Object & Claim type & Scoped claim & Evidence status \\
\hline
P1 & Proposition & Testable Separability Conjecture inside $\Omega$ & Analytic argument; proposed experiment (\S9.1--\S9.2) \\
Path safety / mediation & Proposition & Complete mediation enables path-level admission checks & Analytic argument; cited external evidence; proposed experiment (\S9.3, \S9.5) \\
P15 & Hypothesis & Output-plus-process gates improve discrimination, attribution, or training & Cited external evidence; proposed experiment (\S9.8) \\
P16 & Hypothesis & Pinned summaries support use-case-relative reconstruction & Proposed experiment (\S9.9) \\
P17 & Proposition & Registry schemas/modules remain independent under declared changes & Analytic argument; cited external evidence; proposed experiment (\S9.9) \\
Dry-run / locality & Hypothesis & Planning cost is offset by avoided work or movement & Proposed experiment (\S9.6) \\
Skill lifecycle / training & Design pattern & Train, validate, freeze, monitor, and roll back Skills & Cited external evidence; proposed experiment; future work \\
Higher-order memory & Open question & Structures beyond the contract-relevant IR & Future work \\
\hline
\end{tabular}
\end{table}

\section{Harness Manageability}
\label{sec:manageability}

\subsection{Control Plane and Data Plane}

The control plane owns Skill registration, contract compilation, policy versions, placement rules, and release state. The data plane carries request payloads, executes accepted nodes, moves authorized evidence, and emits traces. The planes interact through references and resolved contracts rather than by copying the full control state into every model prompt.

\begin{figure}[t]
\centering
\caption{Control-plane separation versus control-state leakage. In the managed design, the request receives only activated contracts and opaque version references. In the leaking design, every Skill schema and policy is injected into the request path. Selective activation prevents logical scaling from becoming linear prompt growth and limits exposure of governance metadata.}
\label{fig:control-data}
\end{figure}

Figure 3 highlights a failure mode that is easy to mistake for flexibility. Loading every Skill, application programming interface (API) schema, and policy example into each request leaks control-plane state into the data plane. Besides token cost, the leak creates ambiguous authority: a description visible to the planner may look executable even when its release or policy state has changed. The Harness should first select a small candidate set from registry metadata, then materialize only the accepted versions and constraints needed by the run. Tool Forge \cite{rao2026toolforge} demonstrates that intent-scoped routing can reduce task-flow tool context by 99.2\%, providing direct engineering evidence that selective activation prevents logical scaling from eroding data-plane throughput.

\subsection{Path-Aware Invariants}

Per-Skill approval is necessary but insufficient. Let $G_q=(V_q,E_q)$ be the activated graph for request $q$, with nodes representing Skill operations and edges representing data or control dependencies. The Harness evaluates invariants over the path:

\begin{equation}
\mathrm{Accept}(G_q)=\bigwedge_k \phi_k(G_q,u,B,E,T).
\end{equation}

Examples include separation of duties, data-residency continuity, prevention of read-then-publish paths, maximum cumulative privilege, and mandatory human approval before an irreversible effect. A path can violate an invariant even when each node is locally allowed. Runtime checks must also cover dynamically inserted nodes and retries, since they change the effective graph.

\subsection{Audit, Replay, and Reversibility}

An audit trace should distinguish three questions: what the planner proposed, what the deterministic gates accepted, and what the bound Scaffold actually executed. Conflating these stages makes incident review nearly impossible. At minimum, the trace records request and principal identity, candidate and activated Skills, contract versions, policy decision inputs, graph hash, binding decisions, external calls, evidence artifacts, and final effects.

Replay has levels. \textit{Decision replay} recomputes gate outcomes from recorded inputs. \textit{Plan replay} reruns planning under the recorded model and prompt version, accepting that stochastic output may differ. \textit{Execution replay} re-executes against pinned data or service snapshots when available. The runtime should declare which level it supports rather than promising exact replay across mutable external systems.

Manageability also requires a stop path. Skill releases can be disabled, policy versions rolled back, bindings revoked, and in-flight graphs cancelled at Harness checkpoints. These mechanisms turn a trace from passive observability into operational control.

\section{High Cohesion, Loose Coupling, and Context Partitioning}
\label{sec:cohesion}

\subsection{First Principle: The Agent's Job Is Organizing Context}

A single LLM inference produces output determined by two components: the model's internal knowledge (weights, fixed within a task) and the context (variable input). Since the model weights cannot be changed at runtime, the only variable an agent loop can manipulate is the context. Planning, data retrieval, tool invocation, and memory writes are all, at bottom, decisions about what to place in the context for the next inference step.

This observation has a direct consequence for performance under serving systems that expose prefix-cache reuse (for example, vLLM or SGLang). Cache identity requires more than semantically equivalent prompts: the model and tokenizer versions, system and policy metadata, tool serialization and ordering, and every prefix token must remain fixed. If a single token changes in the middle of the prefix, all subsequent positions must be recomputed. The reusable part is not ``the context that changed slightly'' but ``the context whose prefix is frozen token-by-token with changes only at the tail.'' Heterogeneous tools or goals can therefore reduce cache reuse and introduce irrelevant context, but this is one candidate mechanism for long-horizon degradation rather than a complete explanation.

\subsection{Context Partitioning by Tool-Set and Data Source}

The corresponding design heuristic is: \textit{place version-pinned stable material (system prompt, policy metadata, serialized tool schemas, and fixed memory) at the front in deterministic order; place per-turn task content at the tail.} Under a prefix-caching serving stack, a longer token-identical prefix creates a larger reuse opportunity. Stable semantic content alone is insufficient if serialization, ordering, tokenizer, or model versions differ.

One useful axis for partitioning context is the tool-set. A sub-agent that repeatedly uses the same version-pinned set of tools and Harness contracts can maintain a stable prefix when its prompt renderer is deterministic. Conversely, switching among heterogeneous tool-sets may rewrite that prefix and dilute attention. Stable tool membership therefore helps prefix stability but does not guarantee it; the system must measure token identity and cache hits directly. A related heuristic is to divide sub-agents by minimally overlapping tool-sets when doing so preserves task semantics.

The secondary axis is data. Detail such as raw tables, full document text, and verbose intermediate results should not flow between sub-agents as context. Instead, detail is written to the data subsystem (memory or lake), and only summarized context (conclusions, handles, key constraints) is passed forward. Downstream sub-agents retrieve detail on demand via semantic join, not by carrying it in their context. This reduces inter-agent communication from $O(\text{full context})$ to $O(\text{summary})$, raising orchestration efficiency and keeping each sub-agent's context focused.

\subsection{Four Mechanisms for High Cohesion and Loose Coupling}

The partitioning principles above yield four concrete design mechanisms that span all three runtime layers:

\begin{enumerate}
\item \textbf{Tool-set cohesion (Skills layer).} A sub-agent should use one version-pinned tool-set and deterministic prompt serialization throughout a run. This increases the opportunity for prefix stability and KV-cache reuse, which must still be measured. The anti-pattern is an agent that changes tool schemas, ordering, or policy metadata across successive turns without recording the resulting cache discontinuity.

\item \textbf{Minimally overlapping partitioning (parallelism).} When a task is decomposed into parallel sub-agents, their tool-sets should be as disjoint as task semantics allow. Disjoint tool declarations are only a proxy for a smaller shared surface: compute pools, locks, external services, authorization backends, data dependencies, and scheduler state may still create contention or correlated decisions and must be measured separately.

\item \textbf{Summarized-context handoff (loose coupling).} Sub-agents pass only summarized conclusions and handles to memory, not full context. The receiving sub-agent retrieves detail on demand from the data subsystem. This creates a loosely-coupled sequential relationship: dependencies flow through thin summary interfaces, while detail is absorbed by the memory layer. The data subsystem sits outside the three-layer stack, so this absorption does not introduce coupling into the runtime core.

\item \textbf{Seed-affinity placement (Scaffold layer).} Each type of sub-agent has a derivation seed, an image with pre-warmed tools, dependencies, and cache. Same-seed sub-agents are co-located to reuse image layers, tool processes, and KV-cache hot regions (affinity); different-seed sub-agents are spread across resource domains to reduce contention (anti-affinity). This maximizes physical locality within a type and minimizes cross-type interference.
\end{enumerate}

The parallelism argument of mechanism 2 can be stated as a calibration hypothesis. If a fraction $s$ of a workload remains serial, Amdahl's law caps the speedup at $1/(s+(1-s)/N)$ for $N$ parallel sub-agents. We hypothesize $s=s_0+\lambda\rho+\eta z$, where $\rho$ measures tool-set overlap and $z$ collects independently measured shared-compute, lock, service, authorization, data, and scheduler coupling. We likewise hypothesize $E=E_0-\mu\bar{c}$ for orchestration efficiency $E$ and average handoff cost $\bar{c}$. Neither linear relation is implied by the architecture; coefficients, residual structure, and useful operating ranges must be estimated under Section~9. Tool partitioning raises the parallelism ceiling only if reductions in $\rho$ are not offset by the other terms.

The four mechanisms can be composed: mechanisms 1 and 4 target intra-agent locality, while mechanisms 2 and 3 target inter-agent coupling. Tool-set partitioning is an observable design variable across the Skills, orchestration, and Scaffold-placement layers, but the proposed effects remain hypotheses until token identity, cache reuse, contention, handoff cost, and end-to-end behavior are measured together.

\subsection{Derivation-Closure Orchestration}
\label{sec:derivation-closure}

The four mechanisms describe how a partition should look once chosen. They do not say how the runtime arrives at a partition for a request whose required tools and sources are not known in advance. We close that gap with a rule we call \textit{derivation closure}: a sub-agent is instantiated with a relatively fixed tool-set and data-set, and any reference it raises to a tool or source outside that set is not satisfied in place but becomes the admission condition for deriving a new sub-agent.

\paragraph{Bounding the output range.}
Fixing the tool-set and data-set of a sub-agent bounds its reachable output range, because the operations it may perform and the evidence it may cite are both enumerable at admission time. This is the operational reading of high cohesion: the sub-agent is cohesive not because its prompt is focused but because its contract is closed over a specific pair of sets. A request that would widen either set would also widen the output range, which is precisely the event the Harness must observe rather than absorb silently.

\paragraph{References as derivation conditions.}
When a sub-agent encounters a needed search tool or data source that its contract does not cover, granting it inline would erode both the frozen prefix and the closure property. Instead the reference is recorded as an unmet dependency and triggers derivation: the Harness resolves a fresh contract whose tool-set and data-set cover that dependency, seeds a new sub-agent from the matching template, and binds it as a child. Capability growth within a single request therefore proceeds by adding bounded sub-agents rather than by enlarging any existing one. Each derivation event is an ordinary contract compilation and passes the same deterministic gates of Section~5.3, so the request-level graph remains typed and replayable no matter how the partition was discovered.

\paragraph{The main agent as a satisfaction controller.}
The main agent does not execute task tools. Its function is to summarize the outputs of its sub-agents and propose whether the run may terminate. For each output sub-domain $d$, the contract binds a versioned verifier $V_d$, its evidence inputs, and a threshold $\tau_d$ before execution. The controller performs a semantic join across returned summaries, applies top-$k$ selection to retain the best-supported evidence per field, and reports the \textit{satisfaction ratio} of the declared output sub-domains:

\begin{equation}
\mathrm{sat}(\sigma^{out})=\frac{1}{|\mathrm{dom}(\sigma^{out})|}\sum_{d\in\mathrm{dom}(\sigma^{out})} \mathbb{1}\big[V_d(\mathrm{join}_k(\text{summaries})_d)\geq\tau_d\big].
\end{equation}

When every $V_d$ is deterministic or version-pinned with a preregistered error tolerance, the ratio can participate in a contract-level termination predicate. When a verifier is model-based, mutable, or uncalibrated, the ratio is an estimate and cannot by itself authorize termination; the contract must require an independent deterministic gate or named human approval. Subject to that distinction, the loop continues while the ratio is below its target and some unsatisfied sub-domain still admits an untried derivation. Unsatisfied sub-domains identify what is missing, and unmet references supply derivation candidates, so the controller re-invokes existing sub-agents or derives new ones until the target is certified, the derivation budget is exhausted, or no candidate remains. The last two cases are typed outcomes whose unsatisfied sub-domains and verifier versions appear in the trace.

\begin{figure}[t]
\centering
\caption{Derivation-closure orchestration. Each sub-agent holds an admitted tool-set and data-set; references outside those sets become conditions for deriving a new sub-agent rather than being granted inline. The main agent summarizes returns, performs a semantic join with top-$k$ selection, and reports a satisfaction ratio computed by version-bound per-domain verifiers. Only deterministic or calibrated, version-pinned verifier results may participate directly in termination; other scores remain estimates requiring an independent gate.}
\label{fig:derivation-closure}
\end{figure}

Figure 4 shows the loop. Three properties are testable. First, version-bound verifier results can move termination from an implicit model judgment into an inspectable contract predicate, subject to verifier validity. Second, because derivation is the only way to widen tool or data access, the trace is a complete record of admitted and activated capabilities, not a claim about every capability the task counterfactually needed. Third, the per-sub-domain structure aligns orchestration with the training gate of Section~8.5, allowing the same declared decomposition to be evaluated at training and run time.

\subsection{Contract-Compatible Mechanisms}

The Harness contract can encode which tools a sub-agent may invoke (mechanism 1), declared independence candidates (mechanism 2), summarized handoff schemas (mechanism 3), and seed templates (mechanism 4). That makes the mechanisms implementable and auditable through one boundary, but does not guarantee prefix stability, physical independence, low handoff cost, or favorable placement. A planning-time dry-run can probe closure, estimate token-prefix stability, identify shared resources, and check capacity before committing an external effect; Section~9 tests whether the resulting predictions are useful.

\section{Architecture-Compatible Mechanisms}
\label{sec:consequences}

\subsection{External Data as a Contract Instantiation}

External data should not be a privileged side channel into the agent. It is a Harness-contract instantiation with explicit authentication, authorization, schema, snapshot, provenance, freshness, residency, and evidence obligations. A data Skill may discover or query assets, but the resolved contract determines what can be fetched and what evidence must accompany the result.

\begin{figure}[t]
\centering
\caption{External-data access as a Harness-contract instantiation. Source discovery and query planning remain semantic operations, while credentials, snapshots, provenance, residency, and evidence packaging are resolved and enforced at the contract boundary. Data governance composes with agent execution without embedding infrastructure credentials or source-specific policy in the Skill.}
\label{fig:external-data}
\end{figure}

Figure 5 shows how this boundary preserves both portability and accountability. The same analytical Skill can run against different warehouses or document stores when their adapters satisfy the contract. Conversely, the same data source can support multiple Skills without learning their task semantics. Evidence returned to the model is a typed bundle carrying source and snapshot identity, not an unqualified text fragment.

\paragraph{Hypothetical AI4Science walkthrough.}
Consider a request to screen perovskite candidates for a target band gap near 1.3~eV using recent literature and to produce a cited, reproducible report. This is an architecture walkthrough, not an implemented case study or empirical validation. Its explicit chain is \textit{intent $\rightarrow$ contract compilation $\rightarrow$ path gate $\rightarrow$ binding $\rightarrow$ trace $\rightarrow$ typed failure $\rightarrow$ recoupling}. At \textit{intent}, the request fixes the band-gap target together with recency, citation, and reproducibility constraints. During \textit{contract compilation}, the Harness activates literature-review, materials-screening, and report-generation Skills and closes their typed inputs, outputs, effects, and evidence obligations. The \textit{path gate} checks data authorization, institutional HPC identity, a network allowlist for literature and materials databases, and effect closure before any fetch or computation. \textit{Binding} resolves approved databases, a microVM for scripts, and a compatible GPU or HPC Scaffold without exposing infrastructure credentials to the Skills. The \textit{trace} pins Skill and model versions, literature and materials-data snapshots, binding and scheduler decisions, screening-script identity, and citation evidence. A missing authorization, ungrounded citation, unit anomaly, or absence of compatible HPC is returned as a \textit{typed failure}, rather than repaired through an undeclared path. Finally, \textit{recoupling} is visible if a shared HPC quota, scheduler feedback, or resource-sensitive scientific output makes a Scaffold change alter the admission decision, activated path, authorized effects, binding constraints, postconditions, or scientific result. The example therefore identifies where semantic invariance can fail; it does not establish that the conjecture holds.

\subsection{From Query Plan to Intermediate Relation}
\label{sec:intermediate-relation}

Section~3.4 described the data subsystem at the contract boundary: an off-policy loop pre-summarizes sources, and the Harness consults these summaries to plan access. Here we make the ``plan'' step of that mapping concrete as an explicit, auditable record we call the \textit{Intermediate Relation} (IR), situated between two orthogonal registries.

\paragraph{The Intermediate Relation.}
The \textit{Data Wiki} holds what the data is: per-source semantic summaries produced by the off-policy loop (a natural-language summary, format, retrieval keywords, and cross-source relations). For complex reports whose semantics an agent cannot reliably guess, such as multi-sheet spreadsheets with hidden formulas and drifting definitions, entries are built under human--agent collaboration and tested by \textit{reconstructive reflection}. P16 is deliberately use-case-relative: before summary construction, the evaluator preregisters the target report family, reconstruction fields, and report-structure descriptor. Reports are separated into training, validation, and held-out sets; the reconstruction agent receives no source credentials or original tables; summaries may be revised from training and validation failures only; and the held-out set is inspected once for the final comparison. The resulting reconstruction rate tests whether a version-pinned summary is sufficient for that preregistered use case, not whether it contains every possible meaning of the source. The \textit{Theme Wiki} holds what a verified output looks like: reusable templates for output types (reports, charts, decks, analytical write-ups, data-model programs) registered only after the corresponding output sub-domain has been trained and validated under the lifecycle of Section~8.5. The two registries are designed to be orthogonal: a data summary should not change merely because output demand changes, and an output template should not bind to one source.

\paragraph{Orthogonal Wikis.}
The IR is the sole explicit coupling point between them. For a given theme it records which source summaries are hit by semantic join, packaged with six contextual dimensions:
\begin{equation}
\mathrm{IR}:\ \mathrm{theme} \mapsto \langle \{\Sigma(\mathrm{src}_i)\},\ \mathrm{5W1H} \rangle,
\end{equation}
where \textbf{When} states validity and update cadence of the data for this theme; \textbf{Where} states the access-domain boundary (intranet, extranet, public internet); \textbf{Who} states data ownership and whether this agent or theme is authorized; \textbf{What} states the business semantics of the fields; \textbf{Why} states the integration logic, that is, why these sources belong together and the reason behind the semantic join; and \textbf{How} states the physical access path (database, lakehouse, JSON, or a Harness tool). The 5W1H dimensions do not bypass semantic join; they are metadata that must accompany its hits, because a hit without them may be stale (no When), unauthorized (no Who), or unreachable (no How). With the IR explicit, the contract mapping of Section~3.4 refines to: intent determines a theme, the IR resolves its source set and 5W1H record, the Harness validates timeliness, authorization, and access, and only then issues data fetches.

\paragraph{Decoupling claim and contribution boundary.}
P17 distinguishes three change classes: an entry-content update within an existing registry schema, a registry schema or module update, and cross-registry propagation. The proposition concerns schema/module independence, not immutability of entries or the IR: a compliant one-sided entry change should not force a schema or module change in the other registry. Some source or output changes will legitimately modify entries in both a registry and the IR, and a genuinely new semantic type may require a local schema extension; the falsification question is whether that extension needlessly propagates across the registry boundary. Protocol~9.9 tests all three classes. Concurrent work supports the mechanism while differing in target. Scoped verification of evolving agentic context \cite{hsu2026grace} keeps a persistent instruction as a typed graph and validates each proposed edit only within the local neighborhood of the nodes it touches, reporting a substantial reliability gain over an otherwise identical flat-text baseline. That result is evidence for the general principle that explicit structure can confine a change's blast radius. The IR applies the principle to a different boundary: it structures the coupling between a data registry and an output registry, so the property under test is mutual schema/module independence under one-sided change. This subsection proposes no new retrieval or data-integration algorithm. Consistent with the pre-built versus on-line trade-off measured by Chen et al. \cite{chen2026metadata}, the contribution is the systems-level reorganization of information the data subsystem already holds into an explicit, auditable relation layer.

\subsection{Parallelism, Locality, and Planning-Time Dry-Run}

Once the Harness has a graph, it can expose parallelism that a monolithic prompt obscures. Independent nodes may execute concurrently; dependent nodes remain ordered by graph edges and effect constraints. The Scaffold pool contributes locality and capacity facts. The binding algorithm can then co-locate data-intensive nodes, isolate high-risk effects, and reserve scarce accelerators only for nodes that declare them.

A planning-time dry-run resolves graph structure, candidate bindings, budget envelopes, policy decisions, and expected evidence without committing external effects. It answers operational questions before execution: Which data regions will be crossed? Which nodes can run in parallel? Which irreversible effects require approval? Is the predicted cost within budget? Are all required Scaffold types available?

\begin{figure}[t]
\centering
\caption{Planning-time dry-run and locality-aware execution. The Harness first validates and annotates the graph, then groups independent nodes near required data and binds effects to appropriate isolation zones. Explicit graphs permit concurrency and locality optimization while preserving a reviewable pre-execution decision point.}
\label{fig:dry-run}
\end{figure}

Figure 6 supports a modest claim: graph visibility creates an optimization opportunity; it does not guarantee a speedup. Dry-run adds planning overhead, and locality can conflict with load balance, policy, or accelerator availability. Evaluation must therefore report both avoided execution cost and dry-run cost.

\subsection{Skill-as-Code}

Treating a Skill as a prompt file makes capability growth difficult to review. Skill-as-Code applies a software lifecycle to the complete capability contract: source-controlled declarations, schema linting, effect review, contract tests, sandbox tests, model-version tests, signed release artifacts, staged rollout, monitoring, and rollback. That a Skill deserves the status of a software unit with its own identity, and that identity persistence across update, rollback, and removal is measurable rather than assumed, has since been developed as an ontology in its own right \cite{fan2026skillware}; we take that framing as given. What this subsection adds is narrower and specific to a probabilistic control plane: the point in the lifecycle at which a converged Skill stops being regenerated and becomes frozen code, and the release gate that decides when it may.

\begin{figure}[t]
\centering
\caption{Skill-as-Code lifecycle. A candidate moves through draft, training, validation, staged release, and monitoring. Drift triggers revalidation; operators may freeze or thaw a Skill, roll back a staged or active release, and retire an identity while retaining lifecycle continuity and audit evidence.}
\label{fig:skill-as-code}
\end{figure}

Figure 7 makes validation and reversibility part of the release architecture rather than a documentation convention. The lifecycle is \textit{draft $\rightarrow$ train $\rightarrow$ validate $\rightarrow$ staged release $\rightarrow$ monitor}; observed drift returns the Skill to revalidation, a thaw returns frozen behavior to training, rollback restores a prior staged release, and retirement closes activation while preserving identity history. Operation closure is a validation gate: tests must show that externally visible behavior is expressible through declared effects. Model-version stability is tested at the same boundary, and a release is blocked if contract adherence, path selection, or postcondition satisfaction regresses.

Skill-as-Code also serves as a deterministic anchor in the probabilistic control plane. A Skill's operation space closes once its tool-set is stable, its output format is structured, and its decision branches are enumerable; at that point the execution process can be frozen into code rather than regenerated probabilistically each turn. This removes two sources of uncertainty simultaneously: model sampling (code is a deterministic function) and context organization (the code embeds a stable procedure). This adds a third layer to the probabilistic control plane: probabilistic planning, deterministic execution (frozen code), and deterministic gates. The frozen prefix of Section~7.1 stabilizes the context; the frozen procedure stabilizes the decision within it. Together, they turn the high-cohesion design law into an enforceable property: the stable prefix is frozen for cache locality, and the stable procedure is frozen as code for determinism.

\subsection{Training Before Freezing: A Dual-Subgoal Reward}
\label{sec:training-before-freezing}

Skill-as-Code explains how a stabilized Skill is governed and frozen, but not how the Skill reached stability in the first place. SkillOpt \cite{wang2026skillopt} establishes the missing premise: a Skill is the trainable external state of a frozen agent, edited under optimizer discipline (bounded edits, validation gates, rejected-edit buffers) in text space. Its publicly acknowledged limitation is the reward: a single scalar $r(s)\in[0,1]$ works best when the task has automatic verifiers, while open-ended goals, whose success is subjective and multi-dimensional, require stronger evaluation. A scalar reward is a lossy projection of a natural-language goal: it cannot tell the optimizer whether a failure lies in the outcome or in the process, nor which output sub-domain regressed, so convergence must be forced by repeated human judgment.

We propose to structure the reward along two subgoal dimensions, both derived from artifacts the contract architecture already maintains. The vector-valued gate below is a concrete design hypothesis that instantiates the dual-subgoal intuition, not a claim of logical necessity. The \textit{outcome subgoal} scores the result separately along each sub-domain of the Skill's declared output schema $\sigma^{out}$, with an independent annotated benchmark per sub-domain and a weighted aggregate:
\begin{equation}
r_{\mathrm{out}}(s)=\sum_{d\in\mathrm{dom}(\sigma^{out})} w_d \cdot r_d\big(\mathrm{output}_d,\ \mathrm{benchmark}_d\big).
\end{equation}
The \textit{process subgoal} supplies a dense, immediate signal about how the result was reached: whether the correct tools and data sources were invoked, whether intermediate reasoning satisfied the declared checking criteria, and whether each step gate passed:
\begin{equation}
r_{\mathrm{proc}}(s)=\frac{1}{|\mathrm{steps}|}\sum_k \mathbb{1}\big[\mathrm{step}_k \text{ satisfies the tool/data/logic/gate criteria}\big].
\end{equation}
Each mini-batch computes both subgoal families, and the validation gate becomes vector-valued rather than scalar: an edit is rejected when any protected output sub-domain regresses beyond tolerance or when protected process consistency drops, so the rejected-edit buffer records component-annotated negative feedback instead of an undifferentiated score.

The scalar $r_{\mathrm{out}}(s)$ remains useful as a report statistic but is not the admission object. Let
\begin{equation}
R(s)=\left(\left(r_d(s)\right)_{d\in\mathrm{dom}(\sigma^{out})},r_{\mathrm{proc}}(s)\right).
\end{equation}
Before training, the evaluator fixes the verifier and tolerance $\epsilon_j$ for every component, together with protected components $P$ and target components $T$. A proposed Skill $s'$ is admitted relative to $s$ only when
\begin{equation}
\mathrm{Accept}(s'|s)=
\left[\bigwedge_{j\in P}R_j(s')\geq R_j(s)-\epsilon_j\right]
\land
\left[\bigvee_{j\in T}R_j(s')>R_j(s)+\epsilon_j\right].
\end{equation}
This Pareto-style gate prevents an aggregate gain from hiding a protected-domain regression while still requiring a material improvement on at least one declared target.

The role of the human changes accordingly: instead of judging every iteration, the human confirms the evaluation method once, deciding which benchmark scores each sub-domain and which criteria gate each step, after which reward computation is automatic. Training therefore decomposes naturally: different sub-agents, or per-sub-domain benchmarks, can guide the iteration of a complex Skill separately.

\paragraph{Relation to concurrent work on structured credit assignment.}
Several concurrent systems reach the same conclusion that a single scalar is too coarse, and it is worth stating precisely what remains distinct here. Gated semantic quality-diversity \cite{luo2026gsme} assigns the proposal role to a model and the measurement role to deterministic code, then archives accepted patches by the pathology each one addresses; workflow-localized mechanism learning \cite{lin2026wml} attributes a failure to a specific workflow node and mechanism before selecting the smallest valid edit. Both index credit by an inferred property of the failure, discovered after the fact. Our decomposition instead indexes credit by the sub-domains of the declared output schema $\sigma^{out}$, which exist in the contract before any run occurs. The practical consequence is that the reward dimensions are fixed by the same declaration that the Harness already gates and that Section~7.4 uses to compute the run-time satisfaction ratio, so training-time attribution and run-time completeness share one structure. Whether contract-derived dimensions outperform diagnosis-derived ones is an open question, not something we claim. Independent support for the underlying premise comes from a continual-learning evaluation of harness optimizers, which finds that gains compound only when regression control sits inside the optimization loop \cite{wang2026compound}; a gate that rejects an edit whenever any sub-domain regresses is one way to place it there.

The two stages then compose into one lifecycle. As training converges, subgoals whose evaluation logic and process criteria stop changing have closed operation spaces; they cross the closure threshold of Section~8.4 and enter the freeze-to-code stage. Training and compilation are two phases of the same Skill lifecycle, both governed by the Harness maintenance subsystem. We state the efficiency claim as a falsifiable hypothesis (P15) rather than a result: compared with a scalar gate and an output-vector-only gate, the output-plus-process vector should improve offline discrimination and attribution and improve online convergence without sacrificing protected components at equal rollout budget; Protocol~9.8 specifies separate offline and online comparisons.

\paragraph{Caveat: trustworthiness of process criteria.}
One caveat constrains how the process dimension may be read. A deterministic testbed has shown that harness optimizers can propose guardrails for failure classes that provably never occurred, in 15 of 60 runs when the input merely resembled a familiar rule \cite{wang2026phantom}. Because $r_{\mathrm{proc}}$ depends on judging whether a step satisfied its criteria, a dimension-annotated rejection is only as trustworthy as the criterion that produced it; criteria should therefore be registered artifacts rather than model-generated at scoring time.

\section{Falsifiable Evaluation Protocols}
\label{sec:evaluation}

The present work is architectural and has no controlled implementation results. We therefore specify experiments that can refute its claims. Each protocol separates the intervention from workload, model, policy, and hardware changes. Reported outcomes should include confidence intervals and failure distributions, not only averages.

\begin{figure}[t]
\centering
\caption{Evaluation matrix for the architecture. The P1 rows share a capability-count $\times$ Scaffold-count factorial design with injected recoupling factors. The remaining rows pair a boundary or subsystem hypothesis with a controlled intervention, observations, and a preregistered condition that counts against it.}
\label{fig:evaluation-matrix}
\end{figure}

Figure 8 summarizes the research program. Protocols 9.1 and 9.2 report one randomized capability-count $\times$ Scaffold-count factorial rather than two one-factor demonstrations. Before collection, declare the operating region $\Omega$: workload family; model, tokenizer, Skill, Harness, and policy versions; external-service contracts; compatible Scaffold classes and topologies; and scheduler regime. The unit of randomization is one independent run of a preregistered workload instance under an assigned capability-count and Scaffold configuration. Randomize assignments within workload-instance and seed blocks, and repeat every cell over preregistered model or sampling seeds and independent runs. Before unblinding any outcome estimand, apply the six preregistered mechanism-level acceptance criteria to configuration, binding, information-flow, state-access, resource-channel, and scheduler-channel traces, then freeze and report each pass/fail decision. Outcome values cannot enter those decisions. Report cell estimates and interaction contrasts with confidence intervals and failure distributions.

The shared experiment has three estimands. \textit{Semantic invariance} tests whether preregistered task metrics and the five interface semantics remain equivalent or non-inferior across compatible Scaffold conditions; equivalence or non-inferiority margins must be fixed before collection. \textit{Capacity response} estimates the throughput and latency dose-response over compatible Scaffold count and topology at each capability count. \textit{Recoupling} is the capability-count $\times$ Scaffold-count interaction in those semantic and capacity outcomes, judged against preregistered interaction margins. Separately randomized diagnostic arms inject shared-compute contention, locks, external-service quotas, authorization-backend delay, scheduler feedback, and mutable policy to localize departures. A confidence interval outside a semantic, capacity-response, or interaction margin counts against P1; a non-significant interaction alone is not evidence of independence. The detailed protocols follow.

\subsection{Shared Factorial Experiment: Capability Axis}

\textbf{Hypothesis.} Within $\Omega$, selective Skill activation keeps request-context growth and task interference bounded as unrelated Skills are added, and the capacity-response function for existing workloads remains invariant across capability-count rows.

\textbf{Control.} Execute all randomized factorial cells, treating capability count as the number of registered Skills presented to a fixed retrieval-and-activation mechanism and Scaffold count/topology as the compatible capacity dose. Keep active task Skills fixed and add unrelated Skills in randomized batches. Compare eager prompt loading with registry retrieval followed by typed activation, but analyze each activation strategy as a prespecified stratum rather than changing it across cells. Two confounds require explicit control. First, retrieval competence must be fixed and reported, because a controlled study of progressive disclosure found that its benefit shrinks toward zero when the harness already partitions and retrieves competently, and becomes decisive only at corpus scale \cite{he2026disclosure}. Run both a single-source and a many-source scale. Second, hold activation depth at one level unless depth is the intervention, as the same work reports that a second routing layer never helped and sometimes reduced accuracy.

\textbf{Metrics.} For semantic invariance, record task success, every-requirement satisfaction, plan validity, tool-selection errors, and the admission decision, activated path, authorized effects, binding constraints, and postconditions. For capacity response, record throughput and latency across every Scaffold column. Also report activated-Skill precision and recall, prompt tokens, policy-decision latency, and accidental exposure of control metadata. Aggregate recall is insufficient on its own. A white-box study of Skill execution under long contexts observed requirement coverage remaining above 92\% while task outcomes collapsed, because a small number of dropped requirements can invalidate an otherwise complete artifact \cite{xue2026longcontext}. Report therefore the fraction of runs in which \textit{every} declared requirement was satisfied, alongside aggregate figures, and record which requirements were dropped rather than only how many.

\textbf{Expected signature.} Under selective activation, registry size grows while activated context, semantic outcomes, and capacity-response curves remain within their preregistered margins; eager loading may exhibit increasing context and cross-Skill confusion.

\textbf{Falsification.} The capability-axis claim is weakened if selective activation loses relevant Skills, its retrieval overhead dominates, semantic outcomes leave their equivalence or non-inferiority margins, or capacity-response contrasts vary across capability-count rows beyond the preregistered interaction margin.

\subsection{Shared Factorial Experiment: Capacity and Interaction}

\textbf{Hypothesis.} Within $\Omega$, adding compatible Scaffold instances has a positive capacity response without changing preregistered semantic outcomes or the five interface semantics, and the response does not materially depend on capability count.

\textbf{Control.} Use the same randomized units, blocks, repeated seeds, runs, and factorial cells as Protocol~9.1. Vary only compatible Scaffold count and prespecified topology within each capability-count row; analyze homogeneous and locality-constrained pools as separate strata. Run each recoupling injection as a randomized diagnostic arm, preserving the base cell assignment and recording whether the injection changes a causal condition or leaves $\Omega$.

\textbf{Metrics.} Estimate throughput and p50/p95 latency dose-response curves with confidence intervals, plus queue time, binding failures, isolation failures, cost per completed run, and graph completion rate. Test semantic equivalence or non-inferiority for task metrics and each interface semantic using their preregistered margins. Estimate capability-count $\times$ Scaffold-count interactions for both semantic outcomes and capacity-response parameters with confidence intervals rather than comparing only marginal means.

\textbf{Expected signature.} Queueing delay falls and throughput rises until a shared Harness, data, or external-service bottleneck is reached, while semantic outcomes and interface semantics remain inside their margins and capacity-response curves remain stable across capability-count rows.

\textbf{Falsification.} P1 is rejected in the tested $\Omega$ if all six preregistered mechanism-level condition criteria pass yet a compatible Scaffold intervention leaves a semantic equivalence or non-inferiority interval outside its margin, materially changes an interface semantic, produces a capacity-response departure beyond its margin, or yields a capability-count $\times$ Scaffold-count interaction outside its preregistered bound. Diagnostic-arm failures localize recoupling but do not by themselves show that the base operating region satisfies or violates P1.

\subsection{Harness Bypass Ablation}

\textbf{Hypothesis.} Complete Harness mediation reduces unauthorized effects and improves decision replay.

\textbf{Control.} Execute the same workloads with full mediation and with narrowly instrumented bypasses: direct tool invocation, untyped data fetch, unrecorded retry, and direct Scaffold selection.

\textbf{Metrics.} Count undeclared effects, policy violations, evidence omissions, non-replayable decisions, incident localization time, and false rejections.

\textbf{Expected signature.} Bypass variants may reduce local latency but produce more unexplained or policy-inconsistent runs.

\textbf{Falsification.} If bypass does not measurably affect control or replay under adversarial workloads, the claimed value of the Harness boundary is overstated or the test lacks meaningful effects.

\subsection{Control-Plane Leakage}

\textbf{Hypothesis.} Passing only activated contracts prevents sensitive or irrelevant registry state from entering request context.

\textbf{Control.} Seed the registry with canary metadata, deprecated schemas, and policy text that is never required by the workload. Compare eager loading, retrieved activation, and opaque-reference activation.

\textbf{Metrics.} Track canary reproduction, prompt tokens, deprecated-tool selection, policy-version confusion, and plan accuracy.

\textbf{Falsification.} The claim fails if control metadata appears in outputs or influences plans despite opaque activation, or if opacity prevents the planner from satisfying valid requests.

\subsection{Path-Level Composition Safety}

\textbf{Hypothesis.} Graph-level invariant checks catch hazards that per-Skill approval misses.

\textbf{Control.} Construct Skill pairs and longer paths that are benign individually but violate data-flow, privilege, or effect-ordering constraints when composed. Compare local allowlists with path-aware admission.

\textbf{Metrics.} Measure hazardous-path detection, false positives, dynamic-graph coverage, and decision latency.

\textbf{Falsification.} The claim is weakened if invariants cannot express realistic hazards, if dynamic insertion routinely escapes checks, or if false positives make valid composition impractical.

\subsection{Dry-Run and Locality Economics}

\textbf{Hypothesis.} Planning-time dry-run and locality-aware binding reduce wasted work or data movement enough to justify their overhead for complex graphs.

\textbf{Control.} Hold graph and Scaffold pool fixed. Compare immediate execution with dry-run admission, and locality-oblivious with locality-aware binding.

\textbf{Metrics.} Report dry-run latency, avoided failed-run cost, bytes moved across zones, completion latency, resource utilization, and policy rejection stage.

\textbf{Falsification.} The feature is not justified where dry-run predictions are poorly calibrated, planning cost exceeds avoided work, or locality choices reduce utilization enough to increase total cost.

\subsection{Skill-as-Code Across Model Versions}

\textbf{Hypothesis.} Contract and path tests identify model upgrades that alter Skill behavior before production release.

\textbf{Control.} Freeze Skill source and test corpus while varying planner and executor model versions. Include perturbations of input phrasing, tool errors, and partial data. Separate the component under test from the component that scores it, and grade the first attempt only. Without that separation an agent can silently repair itself mid-run and have the repaired output scored as a pass, which inflates pass rates toward a spurious ceiling \cite{anand2026aeval}. Log every self-correction as an observation in its own right rather than crediting it.

\textbf{Metrics.} Measure schema adherence, declared-effect coverage, path stability, postcondition satisfaction, calibration of abstention, and test flakiness.

\textbf{Falsification.} The lifecycle is insufficient if production-significant changes pass the suite, or if nondeterminism makes the suite too unstable to gate releases.

\subsection{Dual-Subgoal Reward for Skill Training}
\label{sec:eval-dual-subgoal}

\textbf{Hypothesis} (P15). Under a natural-language goal, the output-plus-process vector gate admits more Pareto-valid edits and gives more valid attribution than scalar or output-vector-only gates, improving convergence at equal rollout budget without increasing protected-component regressions.

\textbf{Offline gate comparison.} Construct one fixed proposal bank for a multi-sub-domain task whose outputs combine, for example, numeric answers, textual summaries, and citations. Evaluate every proposal under the same baseline Skill with three gates: (a) one scalar; (b) the output vector $\left(r_d\right)_{d\in\mathrm{dom}(\sigma^{out})}$; and (c) the output-plus-process vector $R(s)$ of Section~8.5. Before scoring, preregister all verifiers, tolerances, protected and target sets, process criteria, decision thresholds, and the two comparator-specific contrasts $(c)-(a)$ and $(c)-(b)$. The offline P15 benefit claim requires both contrasts to meet their preregistered margins; averaging the comparators or passing only one contrast is insufficient. Independent experts or deterministic tests assign blinded oracle labels for Pareto validity, protected-component regression, and process-attribution truth without seeing gate decisions. Grade the executor's first attempt only, keep the grader separate from the executor, and log rather than credit self-corrections \cite{anand2026aeval}. Include counterfactual proposals or episodes with no violation of a selected process criterion as a phantom-violation oracle, following the design that exposed fabricated guardrails in prior harness optimizers \cite{wang2026phantom}.

\textbf{Online training comparison.} Reproduce the SkillOpt \cite{wang2026skillopt} training loop with each of the three gates. Hold model, workload distribution, optimizer prompt, rollout budget, verifier versions, and stopping rules fixed, but allow each loop to generate its own proposals. Run matched training over multiple preregistered random seeds. Preserve first-attempt grading and executor/grader separation. Preregister separate $(c)-(a)$ and $(c)-(b)$ contrasts for convergence, rollout cost, protected-component regressions, and attribution validity; both comparator-specific benefit contrasts must meet their margins. Compare convergence, accepted edits, rollout cost, protected-component regressions, and attribution validity, reporting seed-level distributions and confidence intervals.

\textbf{Metrics.} Offline, report gate discrimination for Pareto-valid versus dominated or regressive proposals, false admission and rejection rates, protected-regression detection, and attribution agreement with blinded oracle truth. Online, report edits or epochs to convergence, final per-domain and process scores, accepted edits, rollout cost, protected regressions, human interventions, and attribution validity.

\textbf{Offline falsification.} The discrimination and attribution part of P15 is weakened if configuration (c) fails a preregistered discrimination, protected-regression-detection, Pareto-valid-admission, or attribution-benefit margin against either (a) or (b), or reports process violations that the blinded phantom-violation oracle disproves.

\textbf{Online falsification.} The convergence and safety part of P15 is weakened if configuration (c) fails a preregistered convergence or rollout-cost benefit margin against either (a) or (b), increases protected-component regressions or degrades attribution validity relative to either comparator, or rejects enough Pareto-valid proposals that matched multi-seed training stalls.

\subsection{Intermediate-Relation Decoupling and Reconstructive Reflection}
\label{sec:eval-ir}

\textbf{Claims} (P16, P17). P16 hypothesizes that a version-pinned Data Wiki summary is sufficient for a preregistered reconstruction use case on a strictly held-out report set. P17 proposes that registry schemas/modules remain independent under compliant one-sided changes while legitimate entry and IR updates remain observable.

\textbf{P16 control.} Preregister the use case, fields, report-structure descriptor, reconstruction metric, and acceptance margins before splitting reports into strict training, validation, and held-out sets. Build and revise summaries using training and validation only, then pin the summary and inspect the held-out set once. On identical held-out tasks, randomize the reconstruction agent to: (a) the pinned summary; (b) token-matched direct reading of source content; (c) metadata-only input; (d) preregistered field ablations of the summary; or (e) shuffled summaries assigned to the wrong source or report family. Deny source credentials and originals to every arm except the bounded direct-reading content supplied to arm (b), and pin the agent, verifier, token budget, and scoring procedure before one-shot held-out use.

\textbf{P16 metrics and contrast.} Report held-out reconstruction accuracy, field-level error, abstention, and confidence intervals for all pairwise preregistered contrasts. Summary-specific sufficiency requires the pinned summary to be equivalent or non-inferior to token-matched direct reading within its margin while exceeding metadata-only, shuffled-summary, and informative field-ablation controls by their positive margins. This joint contrast distinguishes summary structure from a generic context or token benefit.

\textbf{P17 control and artifact-change oracle.} Build a versioned Data Wiki and Theme Wiki over an enterprise-like corpus. Before each intervention, preregister an expected dependency graph for three change classes in each direction: (a) an entry-content update valid under the existing schema; (b) a new semantic type requiring a local registry schema/module update; and (c) a cross-registry use-case change expected to update the IR and named entries. A named \textit{artifact-change oracle} records touched registry entries, schemas or modules, IR records, tests, and generated artifacts, then compares the observed change set with the necessary nodes and edges in the preregistered graph. Reviewers label extra and missing propagation while blinded to the implementation's claimed independence.

\textbf{P17 metrics and falsification.} Measure entry, IR, schema/module, test, and generated-artifact blast radius separately, and classify each observed change as required, missing, or unanticipated relative to the graph. P17 is weakened when a compliant one-sided entry or local schema/module intervention causes unanticipated schema/module propagation into the other registry. Expected entry, IR, test, and generated-artifact changes do not by themselves falsify independence.

\textbf{P16 falsification.} P16 is weakened when the pinned summary fails the joint summary-specific sufficiency contrast, when its one-shot held-out gap versus direct reading exceeds the preregistered margin, or when leakage auditing finds unauthorized source access, held-out-guided summary revision, or an unpinned verifier.

\section{Discussion}
\label{sec:discussion}

\subsection{What the Architecture Does Not Solve}

The Harness contract cannot make an unreliable model correct, an external service deterministic, or a malicious operating system trustworthy. It does not eliminate semantic ambiguity in Skill descriptions, and it does not prove that a policy captures an organization's intent. It can only make selected assumptions explicit and enforce the parts reducible to runtime checks.

The model also introduces a control-plane concentration risk. If the Harness compiler, registry, or policy evaluator is unavailable, execution may stop. If any is compromised, typed contracts can become a false assurance. Implementations therefore need replicated control services, signed artifacts, independent attestation, and a clearly scoped degraded mode. These are engineering requirements, not consequences guaranteed by the abstract model.

The cost of Harness mediation is not negligible and should be budgeted explicitly. Contract compilation, comprising schema validation, graph construction, and invariant evaluation, adds a planning-time overhead that scales with the number of activated Skill nodes and the complexity of path invariants, not with registry size. In practice, this overhead is bounded by the activated graph, not the full Skill set: if a request activates 5 of 200 registered Skills, the gate evaluates 5 nodes and their composition edges, not 200. One data point now bounds this cost from below. Standing-invariant gating in a self-improving runtime has been measured at roughly 0.021~ms per request, a small fraction of inference cost \cite{soni2026gates}. That figure covers invariant evaluation over a bounded model rather than the full contract compilation described here, which also validates schemas, constructs a graph, and resolves bindings, so it is a lower bound rather than a prediction for our setting. We therefore state no estimate of our own: contract compilation latency is a quantity for the dry-run economics protocol (Section~9.6) to measure, and the relevant question is whether it stays within the same order of magnitude as pure invariant checking once schema and binding work are included.

Policy evaluation is cheaper still when rules are compiled to a decision cache indexed by contract hash. The critical design constraint is that this overhead must be amortized: for a request whose LLM inference takes seconds, tens of milliseconds of contract compilation is a small fraction; for a request that returns a single cached lookup, the same overhead dominates and indicates the request should bypass the full pipeline via a fast-path contract.

The degraded mode must be specified, not assumed. When the Harness compiler is unavailable, the runtime should reject new requests rather than fall back to unmediated execution, because unmediated execution violates complete mediation and leaves P1's declared conditions. When the policy evaluator is degraded, the runtime can admit only requests whose contracts are covered by a signed, cached policy decision with a valid freshness window. When the registry is unavailable, already-activated Skill versions remain usable for in-flight requests, but no new Skill versions can be activated. In all cases, the degraded mode narrows the set of admissible operations rather than silently widening them. The key invariant is: degradation reduces capability, never governance.

\subsection{Trade-Offs}

Selective activation can miss a relevant Skill. Strict operation closure can reject useful exploratory behavior. Path checking can be computationally expensive for dynamic graphs. Trace retention may conflict with privacy and data minimization. Replay may be impossible when external systems are mutable. Physical decoupling may be limited when a Skill requires specialized hardware or data residency. The architecture makes these couplings visible; it does not wish them away.

\subsection{Research Questions}

Several questions remain open. What is the smallest contract language that captures effects without becoming a second programming language? Which graph invariants can be checked statically, and which require runtime monitors? How should uncertainty from model-generated planning be represented in a deterministic admission decision? Which interventions truly preserve typed closure, complete mediation, effect non-interference, shared-state isolation, resource invariance, and scheduler independence inside a useful $\Omega$? Which semantic-invariance margins and capacity-response models are scientifically defensible, and which diagnostic injections best localize recoupling? How can traces make those violations attributable while minimizing sensitive content? At what workload complexity does planning-time dry-run become economical? Whether higher-order memory beyond the contract-relevant IR improves these outcomes remains an open question. These questions define the next implementation phase.

\section{Limitations and Threats to Validity}
\label{sec:limitations}

The paper presents a reference architecture, not a production system or empirical result. The model may underrepresent human approvals, long-running state, multi-tenant incentives, and cross-organization trust. The evaluation protocols could favor implementations that expose rich internal telemetry, making comparisons with opaque systems difficult. Synthetic composition hazards may not reflect real incident distributions; real workloads may contain confidential data that cannot be released.

Terminology is another threat. Skill, Harness, and Scaffold are overloaded across agent frameworks. We define them operationally, but mappings to existing systems will be imperfect. In particular, some frameworks combine Harness and Scaffold responsibilities in one service. The model can still be used analytically, but a clean implementation boundary may require refactoring.

Finally, conditional decoupling may hold only within a bounded operating region. A new Skill can require a new hardware class; a new Scaffold can expose timing or locality behavior that changes task quality. The six causal conditions may be incomplete, correlated, or difficult to hold fixed, while deterministic gates and trace identity can expose a violation without preventing it. Equivalence, non-inferiority, and interaction conclusions depend on chosen margins and repeated-run power; a non-significant interaction cannot establish independence. Experiments must record these couplings rather than classify them away. A useful architecture should tell us where semantic invariance or capacity response breaks and which recoupling factor explains the departure.

Two scope cuts deserve explicit mention. The maintenance and self-management subsystem, which covers compaction, persistence, and related system skills governed by the Harness along the physical axis, is treated here only as a boundary requirement and is not designed in this paper; notably, both the Skill training loop of Section~8.5 and the upkeep of IR records are assigned to this subsystem, whose internals remain future work. For the data subsystem, the full four-component formalization (on-policy data-access APIs, governance expressed as data-usage skills, lifecycle management, and the off-policy indexing loop sketched in Section~3.4), together with higher-order memory structures, is developed in companion notes; this paper incorporates only its contract-relevant part: the Intermediate Relation with its 5W1H record and the Data Wiki / Theme Wiki separation of Section~8.2.

\section{Conclusion}
\label{sec:conclusion}

Scalable agentic runtimes have two different growth problems: acquiring more capabilities and executing more governed work. The Skill-Harness-Scaffold model places a typed, auditable contract between task semantics and physical execution. Its central contribution is the Testable Separability Conjecture: inside a declared operating region, typed closure, complete mediation, effect non-interference, shared-state isolation, resource invariance, and scheduler independence are hypothesized to permit capability and capacity interventions to vary independently. Deterministic gates and trace identity make departures detectable, auditable, and attributable; they do not cause independence.

This framing changes the systems question from "how many tools can the agent see?" to "which capability path was activated, under which contract, on which resources, with which evidence and effects?" The paper proposes estimands, not completed experiments: semantic invariance, capacity response, and their capability-count $\times$ Scaffold-count interaction. External-data contracts remain architecture; context locality and planning-time dry-runs remain components of the dry-run/locality hypothesis; the Intermediate Relation's registry-independence claim remains the P17 proposition; and the Skill lifecycle remains a design pattern. As Table~\ref{tab:propositions} classifies these items, none is a logical consequence of P1. The next step is empirical: implement the boundary, calibrate $\Omega$ and the margins, run the factorial and subsystem protocols, and report both where separation holds and which recoupling factors break it.

\bibliographystyle{plain}
\bibliography{references}

\end{document}
